\subsection{EG normalization freedom as a finite counterterm basis (T1S-PT infrastructure)}
\label{app:eg_normalization_freedom_basis}

\AuditTag This appendix refines the strong-closure carrier (Appendix~\ref{app:eg_causal_perturbation_framework}) from a conceptual description into a PT-facing object: an explicit finite-dimensional parameterization of the admissible normalization freedom at a fixed EFT truncation.
The goal is to remove the ambiguity of phrases like ``finite renormalization freedom'' by making the relevant vector spaces and their constraints auditable.

\subsubsection{Admissible local functional space at fixed truncation}
\paragraph{Field content and grading.}
Fix the field content of Appendix~\ref{app:eft_representative_construction} augmented by ghosts/antighosts and auxiliary fields as needed for gauge fixing (Appendix~\ref{app:ward_brst_consistency}).
Let $\mathrm{gh}(\cdot)$ denote ghost number and $\dim(\cdot)$ the engineering dimension.

\paragraph{Admissible counterterm space.}
For a declared truncation $(D_{\max},L_{\max})$, define the admissible counterterm space as
\[
\mathcal{C}_{\le D_{\max}}^{(\mathrm{gh}=0)}
:=\mathrm{span}\{\ \mathcal{O}\ \mid\ \mathcal{O}\text{ local polynomial functional, }\mathrm{gh}(\mathcal{O})=0,\ \dim(\mathcal{O})\le D_{\max}\ \},
\]
modulo total derivatives and algebraic identities (integration by parts, Bianchi identities, etc.).

\AuditTag In EG, the normalization freedom at order $n$ is represented by local distributions supported on the total diagonal.
At fixed truncation, these freedoms reduce to a finite set of local counterterms living in $\mathcal{C}_{\le D_{\max}}^{(\mathrm{gh}=0)}$ and in related graded spaces at other ghost numbers (for ST identities).

\subsubsection{Counterterm basis registry (auditable finite object)}
\paragraph{Basis registry principle.}
\MathTag Choose an explicit basis
\[
\mathcal{B}_{\le D_{\max}}^{(\mathrm{gh}=0)}=\{\mathcal{O}_1,\ldots,\mathcal{O}_K\}
\]
for $\mathcal{C}_{\le D_{\max}}^{(\mathrm{gh}=0)}$ (and similarly for the needed ghost-number sectors).
Then any finite renormalization/counterterm shift at the declared truncation can be written as
\[
\Delta S=\sum_{i=1}^{K} c_i\,\mathcal{O}_i,
\qquad c_i\in\mathbb{R}\ \text{(or a declared coefficient domain).}
\]

\paragraph{Deterministic registry output.}
To keep the basis auditable and reproducible, the manuscript records a deterministic registry summary generated by a script.

\begin{table}[H]
\centering
\small
\setlength{\tabcolsep}{10pt}
\renewcommand{\arraystretch}{1.15}
\begin{tabular}{r r p{9.5cm}}
\toprule
dimension $d$ & ghost \# & admissible counterterm sector (basis principle) \\
\midrule
2 & 0 & Mass-dimension-2 sectors (e.g. gauge-fixing/auxiliary terms if declared) \\
3 & 0 & No CPT-even renormalizable local invariants in the minimal SM gauge sector (registry-specific) \\
4 & 0 & Renormalizable gauge-invariant operators (SM kinetic + Yukawa + Higgs potential) \\
4 & 1 & Ghost-number-1 local functionals controlling ST breaking (candidate anomaly sector) \\
5 & 0 & Dimension-5 EFT operators (e.g. Weinberg operator sector, if included in truncation) \\
6 & 0 & Dimension-6 EFT basis sector (finite operator dictionary at fixed truncation) \\%
\end{tabular}
\caption{EG counterterm-basis registry at fixed truncation. Rows are reproduced by the deterministic script \texttt{scripts/exp\_eg\_counterterm\_basis.py}.}
\label{tab:eg_counterterm_basis_registry}
\end{table}

\subsubsection{PT-facing consequence}
\AuditTag Once the admissible basis is fixed, ``restoring Slavnov--Taylor identities'' becomes a finite-dimensional constraint satisfaction problem:
the ST breaking term lives in a finite-dimensional local space at ghost number $1$, and restoring ST amounts to selecting coefficients $c_i$ (within the admissible freedom) so that the breaking becomes cohomologically trivial.
The algorithmic statement is recorded in Appendix~\ref{app:eg_st_restoration_algorithm}.

